% ************************** Thesis Abstract *****************************
% Use `abstract' as an option in the document class to print only the titlepage and the abstract.
\begin{abstract}
Patch-based diffusion models offer a possible route to strong image priors when training data or memory is limited. This work reproduces and extends the Patch Diffusion Inverse Solver (PaDIS) of Hu et al. \cite{} for computed tomography (CT), using LIDC-IDRI data \cite{} and a fully specified, noise-free fan-beam geometry implemented in LION \cite{}. To separate ambiguities in the original study from adaptations required by the new geometry, reconstruction methods used paper-style and public-compatible reproduction tracks, plus a LION-physics extension. The evaluation covers 8-, 20-, and 60-view CT, limited-angle and $512\times512$ reconstruction, alternative samplers, patch construction methods, and training ablations.

PaDIS produced high-quality reconstructions, including $40.08\,\mathrm{dB}$ PSNR and $0.960$ SSIM at $512\times512$. With LION-physics scaling, PaDIS combined with VE-DDNM achieved the best PSNR at 20 and 60 views, reaching $36.49$ and $41.55\,\mathrm{dB}$. However, the original claim that PaDIS consistently outperforms whole-image diffusion \cite{} was not reproduced. Under VE-DPS, the whole-image model exceeded PaDIS at both 20 and 8 views. Larger patches performed best, positional encoding had little effect on conditioned reconstruction, and using all available slices did not improve PaDIS under a fixed training-time budget. These results show that patch priors are viable and memory-efficient, but their advantage is conditional on the sampler, data regime, and treatment of the forward operator. Reproduction was further limited by incomplete geometry, training, tuning, and implementation details in the original paper and by material differences between its algorithms and released code.
\end{abstract}
